
\documentclass{article}
\usepackage{colortbl}
\usepackage{makecell}
\usepackage{multirow}
\usepackage{supertabular}

\begin{document}

\newcounter{utterance}

\centering \large Interaction Transcript for game `wordle\_withclue', experiment `high\_frequency\_words\_clue\_no\_critic', episode 2 with reflexion{-}Qwen3.5{-}27B.
\vspace{24pt}

{ \footnotesize  \setcounter{utterance}{1}
\setlength{\tabcolsep}{0pt}
\begin{supertabular}{c@{$\;$}|p{.15\linewidth}@{}p{.15\linewidth}p{.15\linewidth}p{.15\linewidth}p{.15\linewidth}p{.15\linewidth}}
    \# & \multicolumn{2}{c}{Player} && \multicolumn{2}{c}{Game Master} \\
    \hline

    \theutterance \stepcounter{utterance}  
    & & & \multicolumn{4}{p{0.6\linewidth}}{
        \cellcolor[rgb]{0.9,0.9,0.9}{
            \makecell[{{p{\linewidth}}}]{
                \texttt{\tiny{[P1$\langle$GM]}}
                \texttt{You have attempted this same task before and did not succeed. These are the plans you wrote after those attempts.} \\
\texttt{Trial \#1: New plan: The previous attempt was aborted because the response violated the strict formatting rules by including a long, rambling explanation that failed to conclude with a valid "guess:" tag and a five{-}letter word. The explanation drifted into a list of rejected words without ever committing to a final guess. Next time, I will strictly adhere to the required format: a concise "explanation:" followed immediately by "guess: {[}word{]}". I will ensure the guess is a valid five{-}letter English word that respects the feedback (excluding red letters 'd', 'a', 'n', 'c' and placing 'e' in a position other than the 5th) and fits the clue. For example, I will select a word like "reels" (if 'r', 'l', 's' are valid) or "beeps" (if 'b', 'p', 's' are valid), ensuring the output contains no extra text or incomplete thoughts.} \\
\\ 
\texttt{Trial \#2: New plan: The previous attempt was aborted because the response failed to provide a complete "guess:" tag and a valid five{-}letter word, instead trailing off with an incomplete list of possibilities. I violated the rule "Do not generate any other text" by including a rambling, unfinished thought process. Next time, I will strictly adhere to the format: a concise "explanation:" followed immediately by "guess: {[}word{]}". I will ensure the guess is a valid five{-}letter English word that respects the feedback (excluding red letters 'd', 'a', 'n', 'c' and placing 'e' in the 5th position as indicated by the green feedback) and fits the clue. For example, I will select a word like "voice" or "scale" and ensure the output ends definitively with the guess, containing no extra text or incomplete sentences.} \\
\\ 
\\ 
\texttt{You are a language wizard who likes to guess words by using the given rules.} \\
\\ 
\texttt{Welcome to Wordle! You have six attempts to guess the target word, a valid English word of five lowercase letters (a{-}z). Please use the tags "explanation:" and "guess:" to provide a concise explanation for each guess.} \\
\\ 
\texttt{To help you make an informed guess, you will receive a clue for the word, such as} \\
\texttt{clue: snowy white.} \\
\\ 
\texttt{Here is an example guess based on the clue:} \\
\texttt{explanation: In the fairy tail Snow White, the girl is killed because she eats a poisoned apple. And the word apple has 5 letters.} \\
\texttt{guess: apple} \\
\\ 
\texttt{After each guess, your answer will be validated, and you will receive feedback indicating which letters are correct (green), which letters are correct but in the wrong position (yellow), and which letters are incorrect (red). This feedback can be useful in determining which letters to include or exclude in your next guess.} \\
\\ 
\texttt{For example, the feedback for "apple" might be:} \\
\texttt{guess\_feedback: a<yellow> p<yellow> p<green> l<yellow> e<red>} \\
\\ 
\texttt{The explanation should contain details about how both the guess\_feedback and clue work together in tandem to formulate a new guess.} \\
\\ 
\texttt{Do not generate any other text.} \\
\\ 
\texttt{Let's begin with your first guess.} \\
\\ 
\texttt{clue: "...called for his fiddlers \_\_\_"} \\
            }
        }
    }
     \\ \\

    \theutterance \stepcounter{utterance}  
    & \multicolumn{4}{p{0.6\linewidth}}{
        \cellcolor[rgb]{0.561,0.737,0.561}{
            \makecell[{{p{\linewidth}}}]{
                \texttt{\tiny{[P1$\rangle$GM]}}
                \texttt{explanation: The clue "...called for his fiddlers \_\_\_" is a famous line from the nursery rhyme "The Grand Old Duke of York," which continues with "to play a tune." However, a more direct fit for a five{-}letter word in a fiddler context, often associated with the phrase "fiddlers green" or similar idioms, might be "green." But looking at the specific phrase "called for his fiddlers," the most common completion in riddles or word games is often "green" (as in "fiddlers green" meaning a place or a state) or perhaps "play." Wait, the specific line is often "The Grand Old Duke of York, he had ten thousand men; He marched them up to the top of the hill, and marched them down again." That doesn't fit. Let's reconsider the clue. "Called for his fiddlers" might refer to a specific song or story. A very common phrase is "fiddlers on the roof," but that's not 5 letters. Perhaps the word is "green" as in "fiddlers green" (a slang term)? Or maybe the clue is referencing "The Fiddler of Dooney"? No. Let's try a common word associated with fiddlers. "Music"? No. "Tune"? No. "Song"? No. "Play"? No. "Dance"? No. "Reel"? A reel is a type of dance music played by fiddlers. "Reel" is 4 letters. "Reels" is 5. "Fiddlers" often play "reels". Let's try "reels".} \\
\texttt{guess: reels} \\
            }
        }
    }
    & & \\ \\

    \theutterance \stepcounter{utterance}  
    & & & \multicolumn{2}{p{0.3\linewidth}}{
        \cellcolor[rgb]{0.9,0.9,0.9}{
            \makecell[{{p{\linewidth}}}]{
                \texttt{\tiny{[GM$|$GM]}}
                \texttt{attempts = 1} \\
\texttt{target\_word = three} \\
\texttt{target\_word\_clue = "...called for his fiddlers \_\_\_"} \\
\texttt{guess = reels} \\
\texttt{guess\_feedback = r<yellow> e<yellow> e<yellow> l<red> s<red>} \\
            }
        }
    }
    & & \\ \\

    \theutterance \stepcounter{utterance}  
    & & & \multicolumn{4}{p{0.6\linewidth}}{
        \cellcolor[rgb]{0.9,0.9,0.9}{
            \makecell[{{p{\linewidth}}}]{
                \texttt{\tiny{[P1$\langle$GM]}}
                \texttt{guess\_feedback: r<yellow> e<yellow> e<yellow> l<red> s<red>} \\
\\ 
\texttt{Provide your response only in this format.} \\
\texttt{explanation: details} \\
\texttt{guess: word} \\
            }
        }
    }
     \\ \\

    \theutterance \stepcounter{utterance}  
    & \multicolumn{4}{p{0.6\linewidth}}{
        \cellcolor[rgb]{0.561,0.737,0.561}{
            \makecell[{{p{\linewidth}}}]{
                \texttt{\tiny{[P1$\rangle$GM]}}
                \texttt{explanation: The feedback indicates that 'r' and 'e' are in the word but not in the positions guessed (1st and 2nd/3rd), while 'l' and 's' are incorrect. The clue "...called for his fiddlers \_\_\_" likely refers to the phrase "fiddlers green" or a specific song title, but given the letters 'r' and 'e' are present, and 'l'/'s' are out, I need a 5{-}letter word with 'r' and 'e' that fits the context. The phrase "fiddlers green" is a common idiom, but "green" has 'g', 'n' which haven't been tested yet. However, looking at the letters available, "reels" was close but had wrong positions for 'r' and 'e' and wrong letters 'l', 's'. If 'r' and 'e' are yellow, they must be in other spots. Maybe the word is "creed"? No. "freed"? No. Let's reconsider the clue. "Called for his fiddlers" might be from a specific story. Wait, the phrase "fiddlers green" is often used. But if 'l' and 's' are red, "reels" is out. What if the word is "green"? 'g', 'r', 'e', 'e', 'n'. 'r' and 'e' are in there. 'l' and 's' are not. This fits the exclusion of 'l' and 's'. The clue might be a misremembered line or a specific riddle where the answer is "green". Let's try "green".} \\
\texttt{guess: green} \\
            }
        }
    }
    & & \\ \\

    \theutterance \stepcounter{utterance}  
    & & & \multicolumn{2}{p{0.3\linewidth}}{
        \cellcolor[rgb]{0.9,0.9,0.9}{
            \makecell[{{p{\linewidth}}}]{
                \texttt{\tiny{[GM$|$GM]}}
                \texttt{attempts = 2} \\
\texttt{target\_word = three} \\
\texttt{target\_word\_clue = "...called for his fiddlers \_\_\_"} \\
\texttt{guess = green} \\
\texttt{guess\_feedback = g<red> r<yellow> e<yellow> e<green> n<red>} \\
            }
        }
    }
    & & \\ \\

    \theutterance \stepcounter{utterance}  
    & & & \multicolumn{4}{p{0.6\linewidth}}{
        \cellcolor[rgb]{0.9,0.9,0.9}{
            \makecell[{{p{\linewidth}}}]{
                \texttt{\tiny{[P1$\langle$GM]}}
                \texttt{guess\_feedback: g<red> r<yellow> e<yellow> e<green> n<red>} \\
\\ 
\texttt{Provide your response only in this format.} \\
\texttt{explanation: details} \\
\texttt{guess: word} \\
            }
        }
    }
     \\ \\

    \theutterance \stepcounter{utterance}  
    & \multicolumn{4}{p{0.6\linewidth}}{
        \cellcolor[rgb]{0.561,0.737,0.561}{
            \makecell[{{p{\linewidth}}}]{
                \texttt{\tiny{[P1$\rangle$GM]}}
                \texttt{explanation: The feedback shows 'e' is correct in the 4th position (green), 'r' is in the word but not in the 2nd position (yellow), and 'g' and 'n' are incorrect (red). The clue "...called for his fiddlers \_\_\_" likely refers to the phrase "fiddlers green" but since 'g' and 'n' are red, that's not it. However, the phrase "fiddlers on the roof" is famous, but "roof" is 4 letters. Let's reconsider the letters: we need a 5{-}letter word with 'r' (not 2nd) and 'e' (4th). The word "creed" has 'e' in 4th? No, c{-}r{-}e{-}e{-}d. 'e' is 3rd and 4th. Wait, "creed" is c{-}r{-}e{-}e{-}d. 4th letter is 'e'. 'r' is 2nd. But 'r' was yellow in position 2 in "green", meaning 'r' is in the word but not in position 2. So 'r' cannot be 2nd. In "reels", 'r' was 1st (yellow), so 'r' is not 1st. So 'r' must be 3rd or 5th. 'e' is 4th. So pattern: \_ \_ r e \_ or \_ \_ \_ e \_ with 'r' in 3 or 5. Let's try "three"? t{-}h{-}r{-}e{-}e. 'r' is 3rd, 'e' is 4th and 5th. 'r' is not 1st or 2nd. 'e' is 4th. This fits. Does "three" fit the clue? "Called for his fiddlers three"? Maybe. Or "fiddlers" play "three" notes? Unlikely. What about "freed"? f{-}r{-}e{-}e{-}d. 'r' is 2nd (bad). "breed"? b{-}r{-}e{-}e{-}d. 'r' is 2nd (bad). "creed"? c{-}r{-}e{-}e{-}d. 'r' is 2nd (bad). "dread"? d{-}r{-}e{-}a{-}d. 'r' is 2nd (bad). "tread"? t{-}r{-}e{-}a{-}d. 'r' is 2nd (bad). "bread"? b{-}r{-}e{-}a{-}d. 'r' is 2nd (bad). "spread"? 6 letters. "shred"? s{-}h{-}r{-}e{-}d. 'r' is 3rd, 'e' is 4th. 's' was red in "reels", so 's' is out. "chred"? No. "phred"? No. "thred"? No. "} \\
            }
        }
    }
    & & \\ \\

    \theutterance \stepcounter{utterance}  
    & & & \multicolumn{2}{p{0.3\linewidth}}{
        \cellcolor[rgb]{0.9,0.9,0.9}{
            \makecell[{{p{\linewidth}}}]{
                \texttt{\tiny{[GM$|$GM]}}
                \texttt{Error: The response should contain the 'guess:' keyword exactly once.} \\
            }
        }
    }
    & & \\ \\

    \theutterance \stepcounter{utterance}  
    & & & \multicolumn{2}{p{0.3\linewidth}}{
        \cellcolor[rgb]{0.9,0.9,0.9}{
            \makecell[{{p{\linewidth}}}]{
                \texttt{\tiny{[GM$|$GM]}}
                \texttt{game\_result = ABORT} \\
            }
        }
    }
    & & \\ \\

    \theutterance \stepcounter{utterance}  
    & & \multicolumn{4}{p{0.6\linewidth}}{
        \cellcolor[rgb]{0.95,0.95,0.95}{
            \makecell[{{p{\linewidth}}}]{
                \texttt{\tiny{[P1$\leftrightarrow$P]}}
                \texttt{Post{-}episode Reflexion (not sent to GM):} \\
\texttt{New plan: The previous attempt was aborted because the response failed to provide a complete "guess:" tag and a valid five{-}letter word, instead trailing off with an incomplete list of possibilities and internal monologue. I violated the rule "Do not generate any other text" by including a rambling, unfinished thought process. Next time, I will strictly adhere to the format: a concise "explanation:" followed immediately by "guess: {[}word{]}". I will ensure the guess is a valid five{-}letter English word that respects the feedback (excluding red letters 'l', 's', 'g', 'n' and placing 'e' in the 4th position as indicated by the green feedback, while placing 'r' in a position other than 1st or 2nd) and fits the clue. For example, I will select a word like "three" or "shred" (if 's' wasn't red, but it was, so "three" is better) and ensure the output ends definitively with the guess, containing no extra text or incomplete sentences.} \\
            }
        }
    }
    & \\ \\

\end{supertabular}
}

\end{document}
